\section{Characteristics and Inferential Resolution}\label{sec:characteristics}

On the focused 408-occupation support, beta's preperiod-weighted correlation is 0.794 with O*NET computer use and 0.589 with remotability.  A static horse race is therefore descriptive: residual variation may be narrow, and a control may also be an AI channel.

The leading computerization measure is the May 2020 O*NET rating ``Interacting With Computers.''  On common support the beta-plus-Webb coefficient is $-0.107$.  Adding computer use makes it $-0.212$; the paired movement is $-0.105$ ($[-0.183,-0.027]$).  Computer use acts as a suppressor, rejecting the simple claim that this control necessarily attenuates the CPS contrast but not isolating AI.  Remotability alone gives $-0.116$ and an undetected $-0.009$ movement.  Rebuilt-treatment checks of wages, education requirements, and routine/manual content likewise do not yield individually detected movements on their declared maximal supports.  A separately repaired pandemic-shortfall exercise uses bounded predictions and a stricter 363-occupation support; total-stock and young-relative conditioning move the pooled coefficient by only 0.0004 and 0.0011 log point.  Regenerating these controls within linked-household draws widens uncertainty but does not resolve a conditioning movement.  These are exploratory diagnostics, not equivalence results.

\begin{table}[!htbp]
\centering
\caption{Computer-use and broad-family conditioning on common support}\label{tab:characteristics}
\resizebox{\textwidth}{!}{%
\begin{tabular}{llrrr}
\toprule
Family-month paths & Computer-use control & Occupations & Q5 by post [95\% CI] & Change from pooled baseline [95\% CI] \\
\midrule
No & No & 455 & $-0.0958\;[-0.1807,-0.0108]$ & --- \\
No & Yes & 455 & $-0.1966\;[-0.3161,-0.0771]$ & $-0.1009\;[-0.1772,-0.0245]$ \\
Yes & No & 455 & $\phantom{-}0.0354\;[-0.0859,0.1567]$ & $\phantom{-}0.1311\;[0.0352,0.2271]$ \\
Yes & Yes & 455 & $-0.0467\;[-0.1821,0.0888]$ & $\phantom{-}0.0491\;[-0.0642,0.1624]$ \\
\bottomrule
\end{tabular}}
\tabnote{All four post-outcome exploratory rows use the same 455 occupations, rebuilt exposure labels, outcome cells, calendar, and 9,999 common occupation multipliers.  The computer-use rating is standardized with preperiod employment weights and interacted with young $\times$ post.  Family-month rows give every SOC2 family its own young-relative calendar-month path.  The final column uses covariance-preserving paired draws.  The four projections are not additive components, and no residual coefficient is a causal AI effect.}
\end{table}


Broad-family paths and a computer-use control need not move the coefficient alike: the former remove all family-common evolution, while the latter partials one correlated dimension.  Neither is a causal technology decomposition.

Industry conditioning is built from occupation--industry--age--month cells rather than dominant-industry labels.  Allowing 13 broad industries their own young-relative post shifts moves the stratified coefficient from $-0.138$ to $-0.099$; the paired 0.039 interval $[0.008,0.070]$ excludes zero.  It does not measure firm AI adoption.  BA and non-BA estimates are $-0.138$ and $-0.123$; their paired difference is $-0.015$ ($[-0.156,0.125]$), with $\mathrm{MDE}_{80}=0.201$.  Exact-age estimates vary but simultaneous paired intervals do not isolate a particular age mechanism.

The industry row changes the cell objective as well as the nuisance set, so it is reported separately from occupation-level characteristic controls.  Education and exact-age estimates use literal common supports and common draws; no subgroup is selected by its point estimate.  The age coefficients are jointly unequal, but each simultaneous age-minus-pooled interval contains zero.

Table~\ref{tab:inference} separates uncertainty targets.  Occupation wild scores allow serial dependence within occupation.  A 22-family sensitivity permits common broad-family shocks.  Household full-refit multipliers preserve repeated observations and co-residence under released weights but are not full CPS design inference.  None covers exposure labels or route allocation.

\begin{table}[!htbp]
\centering
\caption{Inference sensitivities for the core occupational comparison}\label{tab:inference}
\resizebox{\textwidth}{!}{%
\begin{tabular}{lrrr}
\toprule
Procedure and target & Pooled 95\% CI & Conditioned 95\% CI & Paired-movement 95\% CI \\
\midrule
\multicolumn{4}{l}{\textit{Same pooled/SOC2-by-calendar-month targets}} \\
Occupation wild score, SOC2-by-month & $[-0.221,-0.044]$ & $[-0.161,0.117]$ & $[0.011,0.210]$ \\
Twenty-two-family Webb wild score & $[-0.211,-0.053]$ & $[-0.153,0.110]$ & $[0.006,0.215]$ \\
\addlinespace
\multicolumn{4}{l}{\textit{Sampling-oriented companion with regenerated preperiod labels}} \\
Household full refit, 399 draws & $[-0.186,-0.076]$ & $[-0.124,0.066]$ & $[0.035,0.177]$ \\
\bottomrule
\end{tabular}}
\tabnote{All rows report the reconstructed pooled, SOC2-by-calendar-month, and paired-movement coefficients.  The 22-family row permits common shocks within broad occupational families.  The household row uses 399 mean-one linked-household full refits and regenerates preperiod exposure labels; it is conditional on released CPS final weights and is not complete CPS design inference.  The procedures target different perturbations and are not mechanically combined.}
\end{table}


With 22 family clusters, the paired movement remains detected under Rademacher and Webb weights; the Webb interval is $[0.0059,0.2150]$.  The completed finite-sample audit finds small coefficient bias but target-specific interval performance.  In the empirical observed-effect design, occupation-wild coverage is 0.865 for the pooled coefficient, 0.921 for the family-conditioned coefficient, and 0.889 for their paired movement; family clustering is conservative for the pooled and paired targets but covers the family-conditioned target only 0.895 of the time.  In the clean zero-target design, rejection rates for the family-conditioned coefficient are 6.1 percent under occupation weights, 8.9 percent under family weights, and 5.6 percent for the cross-fit oracle.  No one procedure is therefore validated for all targets and declared designs.  In the rebuilt-treatment audit, elapsed-calendar HAC leaves scalar uncertainty similar and none of the 15 five-parameter covariance matrices across three targets and five lags is indefinite.  Earlier indefinite matrices belong to the superseded historical-treatment audit and remain labeled in the archive.
